\begin{enumerate}

\item La suma y diferencia fijan el remolino de entrada.

\[\fitmath{ct1=\mathrm{uu}\cdot \left(1-\mathrm{rea}-\frac{\mathrm{psi}}{2}\right)}\]

\[\fitmath{ct1=(296)\cdot \left(1-(0.55)-\frac{(0.43)}{2}\right)\approx 69.56\,\mathrm{m/s}}\]

\item Se calcula la componente absoluta de salida.

\[\fitmath{ct2=\mathrm{uu}\cdot \left(1-\mathrm{rea}+\frac{\mathrm{psi}}{2}\right)}\]

\[\fitmath{ct2=(296)\cdot \left(1-(0.55)+\frac{(0.43)}{2}\right)\approx 196.84\,\mathrm{m/s}}\]

\item Se resta U y se obtiene el ángulo desde el eje.

\[\fitmath{beta2=\arctan\left(\frac{\mathrm{uu}\cdot \left(-\mathrm{rea}+\frac{\mathrm{psi}}{2}\right)}{\mathrm{cx}}\right)}\]

\[\fitmath{beta2=\arctan\left(\frac{(296)\cdot \left(-(0.55)+\frac{(0.43)}{2}\right)}{(164)}\right)\approx -0.54382\,\mathrm{rad}}\]

\end{enumerate}
